\section{区域不变性}\label{ch10:section-invariance-of-domain}

拓扑学中的重要定理之一便是这个关于``区域不变性''的定理, 因为它表述了欧氏空间的一个内蕴性质. 这一定理由 L. E. J. Brouwer 于 1912 年证得. 定理结论为: 设 $U$ 是 $\mathbb{R}^{n}$ 中的任意开集, $f\colon U\to\mathbb{R}^{n}$ 是任意连续单射, 那么 $f(U)$ 是 $\mathbb{R}^{n}$ 中的开集, 并且 $f$ 的逆映射是连续的. (数学分析中的反函数定理是在增加了假设映射 $f$ 是连续可微并且具有非奇异的 Jacobi 矩阵的条件下得到的.) 我们来证明当 $n=2$ 时的情况.

\begin{lemma}[homotopy extension lemma]\label{ch10:lemma-homotopy-extension}
设 $X$ 是一个空间, $X\times I$ 是正规的, $A$ 是 $X$ 的一个闭子空间, $f\colon A\to Y$ 是连续映射, 其中 $Y$ 是 $\mathbb{R}^{n}$ 的开子空间. 若 $f$ 是零伦的, 则 $f$ 可以扩充为一个连续映射 $g\colon X\to Y$, 并且 $g$ 也是零伦的.
\end{lemma}

\begin{proof}
设 $F\colon A\times I\to Y$ 是 $f$ 与某常值映射之间的同伦. 这时 $F(a,0)=f(a)$, 并且对任意 $a$, $F(a,1)=y_{0}$. 对任意 $x\in X$ 定义 $F(x,1)=y_{0}$, 这样就将 $F$ 扩充到空间 $X\times 1$ 上. 这时 $F$ 是从 $X\times I$ 的闭子空间 $(A\times I)\cup(X\times 1)$ 到 $\mathbb{R}^{n}$ 的连续映射. 由 Tietze 扩充定理, 这一映射可扩充为连续映射 $G\colon X\times I\to\mathbb{R}^{n}$.

映射 $x\mapsto G(x,0)$ 是 $f$ 的一个扩充, 但它是将 $X$ 映射到 $\mathbb{R}^{n}$ 而不是子空间 $Y$. 为了得到所求的映射, 我们作以下操作: 设 $U$ 是 $X\times I$ 的开子集 $U=G^{-1}(Y)$. 则 $U$ 包含 $(A\times I)\cup(X\times 1)$, 见 \cref{ch10:fig-62-1}. 因为 $I$ 是紧致的, 由管道引理可见存在 $X$ 的一个包含 $A$ 的开集 $W$, 使得 $W\times I\subset U$. 由于 $X$ 同胚于 $X\times I$ 的闭子空间 $X\times 0$, 所以 $X$ 是正规的. 因此可以选取连续函数 $\phi\colon X\to[0,1]$ 使得对任意 $x\in A$, 有 $\phi(x)=0$, 对于任意 $x\in X-W$, 有 $\phi(x)=1$. 映射 $x\mapsto x\times\phi(x)$ 将 $X$ 映射到 $X\times I$ 的子空间 $(W\times I)\cup(X\times 1)$ 中, 其中 $(W\times I)\cup(X\times 1)$ 包含在 $U$ 中. 于是连续映射 $g(x)=G(x,\phi(x))$ 便将 $X$ 映射到了 $Y$ 中. 因为对任意 $x\in A$, 有 $\phi(x)=0$, 所以 $g(x)=G(x,0)=f(x)$. 从而 $g$ 就是所求的 $f$ 的扩充, 映射 $H\colon X\times I\to Y$, 定义为
\[
H(x,t)=G(x,(1-t)\phi(x)+t),
\]
便是 $g$ 与一个常值映射间的同伦.
\end{proof}

\par
\begin{figure}[htbp]
\centering
\includegraphics[width=0.7\textwidth]{ch10/fig-62-1.png}
\caption{同伦扩张引理的构造.}\label{ch10:fig-62-1}
\end{figure}

下面的引理是上一节中零伦引理在某些条件下的逆命题.

\begin{lemma}[Borsuk lemma]\label{ch10:lemma-borsuk}
设 $a$ 和 $b$ 是 $S^{2}$ 中的两个点, $A$ 是一个紧致空间, $f\colon A\to S^{2}-a-b$ 是一个连续单射. 如果 $f$ 是零伦的, 那么 $a$ 和 $b$ 属于 $S^{2}-f(A)$ 的同一分支上.
\end{lemma}

\begin{proof}
因为 $A$ 是紧致的, $S^{2}$ 是一个 Hausdorff 空间, 所以 $f(A)$ 是 $S^{2}$ 的紧致子空间, 同胚于 $A$. 因为 $f$ 是零伦的, 所以从 $f(A)$ 到 $S^{2}-a-b$ 的内射也是零伦的. 因此, 我们只要对于 $f$ 是内射的特殊情形证明引理就足够了. 进而, 我们把 $S^{2}$ 换成 $\mathbb{R}^{2}\cup\{\infty\}$, 并且将 $a$ 设为 $0$, $b$ 设为 $\infty$. 这时引理可陈述为:

设 $A$ 是 $\mathbb{R}^{2}-0$ 的紧致子空间, 如果内射 $j\colon A\to\mathbb{R}^{2}-0$ 是零伦的, 那么 $0$ 便属于 $\mathbb{R}^{2}-A$ 的无界分支.

现在我们来证明它. 设 $C$ 是 $\mathbb{R}^{2}-A$ 的包含 $0$ 的分支, 我们假设 $C$ 是有界的, 并且由此推出矛盾. 设 $D$ 是 $\mathbb{R}^{2}-A$ 的其他分支的并, 其中包括无界分支. 那么 $C$ 和 $D$ 是 $\mathbb{R}^{2}$ 的无交开集, 并且 $\mathbb{R}^{2}-A=C\cup D$. 参见 \cref{ch10:fig-62-2}.

我们定义一个连续映射 $h\colon\mathbb{R}^{2}\to\mathbb{R}^{2}-0$ 使得它作用在 $C$ 的补上是恒等映射.

由内射 $j\colon A\to\mathbb{R}^{2}-0$ 入手. 根据假设 $j$ 是零伦的, 由前述引理可见 $j$ 可以扩充为一个从 $C\cup A$ 到 $\mathbb{R}^{2}-0$ 的连续映射 $k$. 这时 $k$ 作用在 $A$ 上时为恒等映射. 再将 $k$ 扩充为映射 $h\colon\mathbb{R}^{2}\to\mathbb{R}^{2}-0$, 其中对任意 $x\in D\cup A$, $h(x)=x$, 从而根据黏结引理 $h$ 是连续的.

现在我们来找出矛盾. 设 $B$ 是 $\mathbb{R}^{2}$ 中的半径为 $M$ 球心在坐标原点的闭球, 其中 $M$ 足够大使得 $B$ 的内部包含 $C\cup A$. (这里要用到 $C$ 是有界的这一性质.) 如果将 $h$ 限制于 $B$ 上, 便获得映射 $g\colon B\to\mathbb{R}^{2}-0$ 使得对任意 $x\in\operatorname{Bd}B$ 有 $g(x)=x$. 将 $g$ 复合从 $\mathbb{R}^{2}-0$ 到 $\operatorname{Bd}B$ 的标准收缩 $x\mapsto Mx/\|x\|$, 就可以获得一个由 $B$ 到 $\operatorname{Bd}B$ 上的收缩. 但这个收缩是不存在的.
\end{proof}

\par
\begin{figure}[htbp]
\centering
\includegraphics[width=0.5\textwidth]{ch10/fig-62-2.png}
\caption{$\mathbb{R}^{2}-A$ 的有界分支 $C$ 与无界部分 $D$.}\label{ch10:fig-62-2}
\end{figure}

\begin{theorem}[invariance of domain]\label{ch10:thm-invariance-of-domain}
如果 $U$ 是 $\mathbb{R}^{2}$ 的一个开子集, $f\colon U\to\mathbb{R}^{2}$ 是连续单射, 那么 $f(U)$ 是 $\mathbb{R}^{2}$ 中的开子集, 并且反函数 $f^{-1}\colon f(U)\to U$ 是连续的.
\end{theorem}

\begin{proof}
像前面所做的一样, 我们用 $S^{2}$ 代替 $\mathbb{R}^{2}$, 证明如果 $U$ 是 $\mathbb{R}^{2}$ 的开子集, $f\colon U\to S^{2}$ 是连续单射, 那么 $f(U)$ 是 $S^{2}$ 的开子集, 并且逆映射是连续的.

\textbf{第一步.} 证明如果 $B$ 是 $\mathbb{R}^{2}$ 中的任何一个包含在 $U$ 中的闭球, 那么 $f(B)$ 不能分割 $S^{2}$.

设 $a$ 和 $b$ 是 $S^{2}-f(B)$ 中的两个点. 因为恒等映射 $i\colon B\to B$ 是零伦的, 所以通过限制 $f$ 得到的映射 $h\colon B\to S^{2}-a-b$ 是零伦的. 根据 Borsuk 引理可见 $a$ 和 $b$ 属于 $S^{2}-h(B)=S^{2}-f(B)$ 的同一分支.

\textbf{第二步.} 证明如果 $B$ 是 $\mathbb{R}^{2}$ 中的任何一个包含在 $U$ 中的闭球, 那么 $f(\operatorname{Int}B)$ 是 $S^{2}$ 中的开集.

因为空间 $C=f(\operatorname{Bd}B)$ 是 $S^{2}$ 中的简单闭曲线, 因此它分割 $S^{2}$. 设 $V$ 是 $S^{2}-C$ 的一个包含连通集合 $f(\operatorname{Int}B)$ 的分支, $W$ 是其他分支的并. 由于 $S^{2}$ 是局部连通的, 所以 $V$ 和 $W$ 是 $S^{2}$ 中的开集. 我们只要证明 $V=f(\operatorname{Int}B)$, 便通过了这一步.

假设 $a$ 是 $V$ 中不在 $f(\operatorname{Int}B)$ 的一个点, 由此推出矛盾. 设 $b$ 是 $W$ 的一个点. 由于集合 $D=f(B)$ 不能分割 $S^{2}$, 所以 $S^{2}-D$ 是包含 $a$ 和 $b$ 的连通集. 再由 $S^{2}-D$ 包含于 $S^{2}-C$ (因为 $D\supset C$) 可得到 $a$ 和 $b$ 属于 $S^{2}-C$ 的同一分支中, 产生矛盾. 参见 \cref{ch10:fig-62-3}.
\end{proof}

\par
\begin{figure}[htbp]
\centering
\includegraphics[width=0.7\textwidth]{ch10/fig-62-3.png}
\caption{区域不变性证明中的 $V,W,C=f(\operatorname{Bd}B)$.}\label{ch10:fig-62-3}
\end{figure}
\begin{exercise}\label{ch10:ex-invariance-of-domain-1}
举例说明当 $f$ 不是单射时 Borsuk 引理的结论可能不成立.
\end{exercise}

\begin{exercise}\label{ch10:ex-invariance-of-domain-2}
设 $A$ 是 $S^{2}$ 的紧致可缩子空间. 证明 $A$ 不能分割 $S^{2}$.
\end{exercise}

\begin{exercise}\label{ch10:ex-invariance-of-domain-3}
设 $X$ 是一个空间, 使得 $X\times I$ 是正规的, $A$ 是 $X$ 的闭子空间, $f\colon A\to Y$ 是连续映射, 其中 $Y$ 是 $\mathbb{R}^{n}$ 的一个开子空间. 如果 $f$ 同伦于一个可以扩充为连续映射 $h\colon X\to Y$ 的映射, 那么 $f$ 自身可扩充为一个连续映射 $g\colon X\to Y$, 使得 $g\simeq h$.
\end{exercise}

\begin{exercise}\label{ch10:ex-invariance-of-domain-4}
设 $C$ 是 $\mathbb{R}^{2}-0$ 中的一条简单闭曲线, $j\colon C\to\mathbb{R}^{2}-0$ 是一个内射. 证明: 若 $0$ 属于 $\mathbb{R}^{2}-C$ 的一个无界分支, 则 $j_{*}$ 是平凡的, 否则是非平凡的. (事实上, 在第二种情况中 $j_{*}$ 是一个同构, 我们将在第 65 节中证明这一点.)
\end{exercise}

\begin{exercise}\label{ch10:ex-invariance-of-domain-5}
设 $U$ 是 $\mathbb{R}^{2}$ 中的一个单连通开集. 如果 $C$ 是包含在 $U$ 中的一条简单闭曲线, 那么 $\mathbb{R}^{2}-C$ 的每一个有界分支都包含于 $U$.

(实际上这一情况刻画了 $\mathbb{R}^{2}$ 中的单连通开集. 参见 $[\mathrm{RW}]$. 当然, 空间 $\mathbb{R}^{2}-C$ 中只包含一个有界分支, 我们将在下一节中证明这一点.)
\end{exercise}

\begin{exercise}\label{ch10:ex-invariance-of-domain-6}
假设你已知不存在从 $B^{n}$ 到 $S^{n-1}$ 的收缩.

\begin{enumerate}[(a)]
\item 证明对于 $S^{n}$, Borsuk 引理成立.
\item 证明 $S^{n}$ 中没有紧致可缩子空间分割 $S^{n}$.
\item 假设你还知道 $S^{n}$ 中的任何一个同胚于 $S^{n-1}$ 的子空间都会分割 $S^{n}$, 证明区域不变性定理在 $n$ 维空间中成立.
\end{enumerate}
\end{exercise}
